\section*{Wstęp i konwencje}
\addcontentsline{toc}{section}{Wstęp i konwencje}
W sekcjach dotyczących jednej zmiennej $[a,b]\subset\R$ jest ustalonym przedziałem zwartym, $a<b$, a $\lambda$ oznacza miarę Lebesgue'a. Zakładamy podstawy teorii miary, algebry liniowej oraz analizy funkcjonalnej. Podstawy grup topologicznych i lokalnej zwartości wprowadzamy przed konstrukcją miary Haara.

\textbf{Co porównujemy.} Całkowalność zależy od funkcji, dziedziny, miary lub funkcji całkującej oraz przestrzeni wartości. Zmiana któregokolwiek z tych elementów może zmienić odpowiedź. O ile nie zaznaczamy inaczej, całki przyjmują wartości skończone. Dla nieujemnych funkcji mierzalnych dopuszczamy również $+\infty$. O ile nie zaznaczamy inaczej, funkcje są rzeczywiste. Dla funkcji zespolonej $f=u+iv$ przyjmujemy, że jest całkowalna w danej teorii, jeżeli całkowalne są obie funkcje rzeczywiste $u,v$; wówczas jej całka jest sumą całki z $u$ oraz $i$ razy całki z $v$. Symbolem $\chi_A$ oznaczamy funkcję charakterystyczną zbioru $A$, równą $1$ na $A$ i $0$ poza $A$.

Dla przestrzeni z miarą $(X,\mathcal A,\mu)$ przez $\mathcal L^1(\mu)$ oznaczamy zbiór konkretnych funkcji $\mathcal A$-mierzalnych $f:X\to\R$ o skończonych wartościach i $\int |f|\,\mathrm d\mu<\infty$, a przez $L^1(\mu)$ jego przestrzeń ilorazową względem równości prawie wszędzie. Jest to istotne przy porównywaniu Riemanna i Lebesgue'a: $\chi_{\Q\cap[0,1]}$ nie jest całkowalna w sensie Riemanna (zob. przykład funkcji Dirichleta w sekcji o całce Darboux oraz twierdzenie o równoważności całek Darboux i Riemanna), chociaż reprezentuje w $L^1(\lambda)$ tę samą klasę co funkcja zerowa: zbiór $\Q\cap[0,1]$ jest przeliczalny, a każdy jego singleton ma miarę Lebesgue'a równą zeru, więc z przeliczalnej podaddytywności miary cały zbiór ma miarę zero.

\textbf{Terminologia.} \emph{Integrand} to funkcja podcałkowa; \emph{integrator} to funkcja całkująca (w teorii Stieltjesa). \emph{Antiderivative} oznacza funkcję pierwotną, zaś \emph{primitive} w teorii całki HK oznacza funkcję pierwotną w sensie tej teorii, tj. funkcję postaci $F(x)=C+\int_a^x f(t)\,\mathrm dt$ dla funkcji $f$ całkowalnej w sensie HK i stałej $C\in\R$. Jej ciągłość oraz równość $F'=f$ prawie wszędzie są własnościami tej całki, a nie warunkami definiującymi funkcję pierwotną. \emph{Gauge} oznacza cechowanie, \emph{tagged partition} --- podział znaczony, \emph{refinement} --- zagęszczenie. \emph{Bounded variation} to wahanie ograniczone, a \emph{regulated function} --- funkcja regulowana (mająca skończone granice jednostronne). \emph{Finite-valued} oznacza funkcję o skończonych wartościach, niekoniecznie ograniczoną.

Podziałem przedziału $[a,b]$ nazywamy zbiór $P=\{a=x_0<x_1<\dots<x_n=b\}$. Przyjmujemy $n\ge1$; jego podprzedziałami są domknięte przedziały $[x_{i-1},x_i]$ dla $1\le i\le n$. Ich wnętrza są parami rozłączne, a sąsiednie podprzedziały mają wspólny koniec. Średnicą podziału nazywamy $\mesh P=\max_{1\le i\le n}(x_i-x_{i-1})$. Po wyborze znaczników $t_i\in[x_{i-1},x_i]$ powstaje podział znaczony $\mathcal P=\{(t_i,[x_{i-1},x_i])\}_{i=1}^n$, o sumie
\[
S(f,\mathcal P)=\sum_{i=1}^n f(t_i)(x_i-x_{i-1}).
\]
W sumach Stieltjesa długość przedziału zastępujemy przyrostem funkcji całkującej.

